% TOP_PROBLEM: {"id": 3007, "title": "Kirby Problem 4.131", "category_id": 7, "category": {"id": 7, "name": "topology", "display_name": "Topology", "description": "Properties preserved under continuous deformations.", "slug": "topology", "order_index": 7, "created_at": "2026-07-31T15:26:25.671Z"}, "status": "open"}
% FIRST_POSTED: null
% ENTRY_KIND: statement_only
% Imported from: batch14.tex; UnsolvedMath ID 3007
\documentclass[11pt]{article}
\usepackage[margin=25mm]{geometry}
\usepackage{fontspec}
\setmainfont{Latin Modern Roman}
\usepackage{amsmath,amssymb,mathtools}
\usepackage{unicode-math}
\setmathfont{Latin Modern Math}
\usepackage{xurl,hyperref}
\hypersetup{colorlinks=true,urlcolor=blue}
\setcounter{secnumdepth}{0}
\setlength{\parindent}{0pt}
\setlength{\parskip}{5pt}
\setlength{\emergencystretch}{3em}
\newcommand{\sref}[2]{\hyperlink{source:#1:#2}{[S#2]}}
\newcommand{\cataloguescope}[1]{\paragraph{Recorded target.}#1}
\begin{document}
% UnsolvedMath ID 3007; source catalogue code KP-4.131
\setcounter{section}{3006}
\section{Kirby Problem 4.131}\label{3007}
{\small\sffamily UnsolvedMath ID 3007 · Topology}
\subsection{Definitions and mathematical statement}

\paragraph{Shared notation.}$\mathbb N=\{1,2,\ldots\}$, and $\mathbb Z,\mathbb Q,\mathbb R,\mathbb C$ denote the integers, rationals, reals and complex numbers. Any other conventions are specified locally.


A Lipschitz structure is an atlas with locally bi-Lipschitz transition maps; compatible smooth structures are those whose smooth charts give this Lipschitz structure. Uniqueness of a smoothing is up to diffeomorphism. For a closed oriented topological four-manifold $X$, put $b_2(X)=\dim_{\mathbb R}H^2(X;\mathbb R)$ and let $\sigma(X)$ be its intersection-form signature. Indefinite means that form has both positive and negative directions. Spin means that the second Stiefel--Whitney class of its oriented topological tangent microbundle vanishes, $w_2(X)=0$. A Lipschitz embedded manifold with boundary in $\mathbb R^4$ is locally bi-Lipschitz equivalent to the standard whole-space or half-space model, as appropriate.
\paragraph{Statement recorded in the supplied source.}
Does every Lipschitz four-manifold admit a compatible smooth structure, and is that structure unique up to diffeomorphism? In addition:
\begin{enumerate}
\item[(a)] Is there a closed spin indefinite topological four-manifold with a Lipschitz structure that violates
$b_2(X)\geq\tfrac54|\sigma(X)|+2$, or that violates
$b_2(X)\geq\tfrac54|\sigma(X)|+4$ when $X$ is not one of $S^4,S^2\times S^2,$ or K3?
\item[(b)] If $X\subset\mathbb R^4$ is a Lipschitz embedded four-manifold with $\partial X\cong S^3$, does $X$ admit a unique smooth structure?
\end{enumerate} \sref{3007}{2}
\subsection{Short English statement}
Are Lipschitz four-manifolds smoothable uniquely, and can a Lipschitz spin manifold evade the smooth spin inequalities; what happens for Lipschitz domains with three-sphere boundary?
\subsection{Sources}
\begin{description}
\item[\hypertarget{source:3007:1}{S1}] ULAM.AI and the UnsolvedMath Contributors, \emph{UnsolvedMath: A Curated Collection of Open Mathematics Problems}, record 3007 (KP-4.131). CC BY 4.0. \url{https://huggingface.co/datasets/ulamai/UnsolvedMath}.
\item[\hypertarget{source:3007:2}{S2}] R. İnanç Baykur, Robion C. Kirby and Daniel Ruberman (eds.), \emph{K3: A New Problem List in Low-Dimensional Topology}, Mathematical Surveys and Monographs 295, American Mathematical Society (2026), Problem 4.131, author preliminary version. \url{https://math.berkeley.edu/sites/default/files/surv-295-ruberman-watermarked-author-pdf.pdf}.
\end{description}
\label{end:3007}
\end{document}
